\subsection{Minimal Polynomial Foundations}

\begin{theorem} \label{thm:operator-minimal-polynomial-existence} \hlt{Existence of the Minimal Polynomial}\\
Let $\mathcal{V}\neq\{\mathbf{0}\}$ be finite-dimensional over $\mathbb{F}$ and let $T\in\LL(\mathcal{V})$. Then:
\begin{enumerate}[label=(\roman*)]
\setlength{\itemsep}{0pt}
\item $T$ has a monic annihilator of positive least degree;
\item the annihilator in (i) is unique.
\end{enumerate}
\begin{proof}
\begin{enumerate}[label=(\roman*)]
\setlength{\itemsep}{0pt}
\item By Proposition \ref{prop:finite-dimensional-operators-have-annihilating-polynomials}, choose a nonconstant annihilator $h$ and replace it by its monic normalization.\\
Let $\mathcal{M}_T$ be the monic annihilators of $T$ and set $\mathcal{D}_T\defeq\{\deg p:p\in\mathcal{M}_T\}$; then $\varnothing\neq\mathcal{D}_T\subseteq\mathbb{N}$.\\
By well-ordering, set $d\defeq\min\mathcal{D}_T$ and choose a monic annihilator $p$ with $\deg p=d$.\\
A nonzero constant $c$ cannot annihilate $T$ because $c(T)=cI_{\mathcal{V}}\neq0$ on $\mathcal{V}\neq\{\mathbf{0}\}$. Thus $d\geq1$.
\item Let $p$ and $q$ be monic annihilators of degree $d$. If $p\neq q$, then $(p-q)(T)=p(T)-q(T)=0$ by Definition \ref{def:polynomial-applied-to-operator}.\\
Their leading terms cancel, so $\deg(p-q)<d$. Normalizing the nonzero polynomial $p-q$ therefore produces a monic annihilator of degree less than $d$, contrary to the definition of $d$. Hence $p=q$.
\end{enumerate}
\hlt{Structural Commentary}: Well-ordering selects a least normalized annihilator; leading-term cancellation forces uniqueness.\\
\hlt{Key Insight}: Two minimal monic annihilators would differ by a smaller annihilator.
\end{proof}
\end{theorem}

\begin{definition} \label{def:operator-minimal-polynomial} \hlt{Minimal Polynomial of an Operator}\\
Let $\mathcal{V}\neq\{\mathbf{0}\}$ be finite-dimensional over $\mathbb{F}$ and let $T\in\LL(\mathcal{V})$. The \hlt{minimal polynomial} $\mu_T$ is the unique monic polynomial of least degree satisfying $\mu_T(T)=0$; it exists by Theorem \ref{thm:operator-minimal-polynomial-existence}.
\end{definition}

\begin{definition} \label{def:minimal-polynomial-matrix} \hlt{Minimal Polynomial of a Matrix}\\
Let $n\geq1$, let $A\in\operatorname{Mat}(n,n,\mathbb{F})$, and let $T_A\in\LL(\mathbb{F}^n)$ be induced by $A$ as in Definition \ref{def:linear-map-induced-by-matrix}. The \hlt{minimal polynomial} of $A$ is $\mu_A\defeq\mu_{T_A}$.\\
For $j\in\mathbb{Z}_{\geq0}$, set $P_j\defeq((\forall\,\mathbf{v}\in\mathbb{F}^n)\,T_A^j\mathbf{v}=A^j\mathbf{v})$. Then $P_0$ holds by Definitions \ref{def:operator-powers} and \ref{def:linear-map-induced-by-matrix}.\\
If $P_j$ holds, then $(\forall\,\mathbf{v}\in\mathbb{F}^n)\,T_A^{j+1}\mathbf{v}=T_A(A^j\mathbf{v})=A^{j+1}\mathbf{v}$ by Definition \ref{def:linear-map-induced-by-matrix}, so $P_{j+1}$ holds. Hence $(\forall\,j\in\mathbb{Z}_{\geq0})\,P_j$ by Theorem \ref{thm:mathematical-induction}.\\
If $p(z)=\sum\limits_{j=0}^{d}c_jz^j\in\mathbb{F}[z]$, then $(\forall\,\mathbf{v}\in\mathbb{F}^n)\,p(T_A)\mathbf{v}=\sum\limits_{j=0}^{d}c_jA^j\mathbf{v}=p(A)\mathbf{v}$ by Definitions \ref{def:polynomial-applied-to-operator} and \ref{def:polynomial-evaluation-on-square-matrix}.\\
Thus $\mu_A$ is equivalently the unique monic polynomial of least positive degree satisfying $\mu_A(A)=0$.
\end{definition}

\begin{theorem} \label{thm:operator-annihilator-iff-multiple-minimal-polynomial} \hlt{Annihilators Are Multiples of $\mu_T$}\\
Let $\mathcal{V}\neq\{\mathbf{0}\}$ be finite-dimensional over $\mathbb{F}$, let $T\in\LL(\mathcal{V})$, and let $f\in\mathbb{F}[z]$ be nonzero. Then $f(T)=0$ if and only if $\mu_T$ divides $f$.
\begin{proof}
\begin{enumerate}[label=\roman*.]
\setlength{\itemsep}{0pt}
\item[$(\Rightarrow)$] Suppose $f(T)=0$ and set $d\defeq\deg\mu_T$. Starting from $r_0\defeq f$ and $s_0\defeq0$, suppose $\deg r_k\geq d$, set $e_k\defeq\deg r_k-d$, and let $a_k\in\mathbb{F}^{\times}$ be the leading coefficient of $r_k$. Set $r_{k+1}\defeq r_k-a_kz^{e_k}\mu_T$ and $s_{k+1}\defeq s_k+a_kz^{e_k}$.\\
The identity $f=s_k\mu_T+r_k$ holds for $k=0$ and is preserved by the recurrence.\\
Monicity of $\mu_T$ implies $r_{k+1}=0$ or $\deg r_{k+1}<\deg r_k$. Degree descent terminates with $f=s\mu_T+r$, where $r=0$ or $\deg r<d$.\\
By Definition \ref{def:polynomial-applied-to-operator} and Proposition \ref{prop:operator-polynomial-products}(i), $r(T)=f(T)-s(T)\mu_T(T)=0$. If $r\neq0$, its monic normalization is an annihilator of degree less than $d$, which is impossible by Definition \ref{def:operator-minimal-polynomial}. Thus $r=0$ and $\mu_T$ divides $f$.
\item[$(\Leftarrow)$] If $f=\mu_Ts$, then $f(T)=\mu_T(T)s(T)=0$ by Proposition \ref{prop:operator-polynomial-products}(i).
\end{enumerate}
\hlt{Structural Commentary}: Leading-term cancellation constructs a remainder; minimality eliminates every nonzero remainder.\\
\hlt{Key Insight}: A nonzero remainder would normalize below $\deg\mu_T$.
\end{proof}
\end{theorem}

\begin{theorem} \label{thm:similarity-preserves-minimal-polynomial} \hlt{Similarity Preserves the Minimal Polynomial}\\
Let $n\geq1$ and let $A,B\in\operatorname{Mat}(n,n,\mathbb{F})$ be similar. Then:
\begin{enumerate}[label=(\roman*)]
\setlength{\itemsep}{0pt}
\item $p_A=p_B$;
\item $\mu_A=\mu_B$.
\end{enumerate}
\begin{proof}
Choose $S\in\operatorname{Mat}(n,n,\mathbb{F})^{\times}$ satisfying $B=SAS^{-1}$ by Definition \ref{def:similar-matrices}.
\begin{enumerate}[label=(\roman*)]
\setlength{\itemsep}{0pt}
\item $p_A=p_B$ by Theorem \ref{thm:similarity-preserves-characteristic-polynomial}(i).
\item Set $P_j\defeq(B^j=SA^jS^{-1})$ for $j\in\mathbb{Z}_{\geq0}$. The identity $P_0$ holds because $B^0=I_n=SI_nS^{-1}=SA^0S^{-1}$. If $P_j$ holds, then $B^{j+1}=B^jB=SA^jS^{-1}SAS^{-1}=SA^{j+1}S^{-1}$, so $P_{j+1}$ holds. Hence $(\forall\,j\in\mathbb{Z}_{\geq0})\,P_j$ by Theorem \ref{thm:mathematical-induction}.\\
By Definition \ref{def:polynomial-evaluation-on-square-matrix}, $(\forall\,f\in\mathbb{F}[z])\,f(B)=Sf(A)S^{-1}$. Therefore $f(B)=0\Longleftrightarrow f(A)=0$, so $A$ and $B$ have the same monic least-degree annihilator by Definition \ref{def:minimal-polynomial-matrix}.
\end{enumerate}
\hlt{Structural Commentary}: Power induction transports characteristic and minimal polynomial data across one conjugation.\\
\hlt{Key Insight}: Conjugation preserves every polynomial equation.
\end{proof}
\end{theorem}
